\chapter{Answers to the self-check questions}
\label{app:answers}

Try each question first. These answers are short on purpose; the chapter gives the detail.

\section*{Chapter \ref{ch:bigpicture}: the big picture}
\begin{enumerate}
  \item 01: five sources describe each place. 02: Earth Engine makes cloud-free composites and exports 15 files per patch; image processing makes the 512-pixel chips and arrays. 03: each patch gets an image pair, an automatic label and an evidence card. 04: the 366 records are frozen and split 256/55/55. 05: three evaluation tracks on the same 55 test patches.
  \item Phase~B is a time (free-form explanations, 2--3 October); Track~B is an experiment (the class LoRA, run in Phase~A).
  \item The test set has 19 increase patches out of 55, so always answering ``increase'' scores 34.5\,\%; that is the score a useless model gets.
  \item Qwen2-VL-2B and LLaVA-Interleave-0.5B.
  \item What the full card's summary line contains; what happens without it (Qwen 56.4\,\% $\rightarrow$ 32.7\,\%); what each model received; the size of the fine-tunes; what the labels are.
\end{enumerate}

\section*{Chapter \ref{ch:stage01}: study design and sources}
\begin{enumerate}
  \item To remove seasonal change: comparing the same months one year apart leaves mostly year-to-year change.
  \item $(B8 - B4)/(B8 + B4)$. Healthy leaves absorb red light and reflect near-infrared strongly, so the numerator is large.
  \item The T0 window straddles ESA's processing change of 25 January 2022, which added an offset; the harmonized collection removes it so the two years are comparable.
  \item VV (vertical in, vertical out) mostly reflects surface scattering and soil moisture; VH (cross-polarised) comes mostly from volume scattering in the canopy, so VH is more sensitive to vegetation.
  \item The drop is the same almost everywhere and 2023 was much drier; drying soil and leaves lower backscatter. Forest-loss alerts need much larger drops in dense time series.
  \item Rainfall comes from 5.5\,km pixels and fell equally for every class; it describes the year, not the patch.
\end{enumerate}

\section*{Chapter \ref{ch:stage02}: Earth Engine and image processing}
\begin{enumerate}
  \item The export used degrees (EPSG:4326); a 640\,m square is about 70 pixels north--south and 65 east--west at this latitude.
  \item It is identical in both years for every patch, which real cloud masks would not be; it probably measures edge pixels.
  \item It multiplies every average by the share of real pixels: about 0.93 for a typical chip and far less for edge patches, so the card's values are biased towards zero.
  \item The threshold that best splits a set of values into two groups (maximum between-group variance).
  \item A policy removed boxes from every patch without a human \texttt{decline} label, and 347 patches had no human label.
\end{enumerate}

\section*{Chapter \ref{ch:stage03}: labels and cards}
\begin{enumerate}
  \item The 15\,\% of the pool with the largest NDVI drop is decline, the 15\,\% with the largest gain is increase, and the rest is stable.
  \item The whole region lost greenness, so the top 15\,\% starts at $\Delta = +0.003$; with the second batch's recomputed cut-offs, an increase patch can have $\Delta$ down to $-0.001$.
  \item Version~1: template text (Ollama failed), only in \file{dataset_old.jsonl}. Version~2: rule-based text computed from the rasters, in the final \file{dataset.jsonl}.
  \item The card's last line, which states \texttt{candidate\_class}; it means a model given the full card has the answer in its input, so Track~A measures use of the complete card.
  \item They use different thresholds on different numbers: percentiles of the unbiased Earth Engine mean for the label; fixed cut-offs on the padded canvas mean for the trend.
  \item Zero padding is included in the canvas average.
\end{enumerate}

\section*{Chapter \ref{ch:stage04}: dataset and split}
\begin{enumerate}
  \item The development set is for making choices and may be used repeatedly; the test set is for the final report only.
  \item 19 of 55 test patches are \texttt{increase}.
  \item Label information in the input (applies: the summary line); spatial leakage between neighbouring patches (applies: random split); geographic memorisation (prevented: coordinates removed).
  \item Zero-shot models learn nothing from the training patches; a fine-tuned model could memorise a neighbourhood.
  \item Content: none of them searched the prompts for the label (\file{check_no_leak.py} was added later).
\end{enumerate}

\section*{Chapter \ref{ch:models}: models}
\begin{enumerate}
  \item A vision encoder, a projector, and a language model.
  \item A SigLIP image encoder and a Qwen1.5-0.5B-Chat language model.
  \item Qwen and LLaVA saw both images and the full card; PaliGemma saw both images and one NDVI line; BLIP-2 saw one stitched image and a short prompt.
  \item It repeated its prompt, and the parser read ``decline'' in the echo for every patch.
  \item $W_0$ is the frozen weight matrix, $B$ and $A$ the trainable thin matrices, $r$ the rank (8) and $\alpha$ the scale (16). A rank-8 adapter adds $2 \times 1536 \times 8 = 24{,}576$ weights.
  \item So that the loss is computed only on the answer; otherwise the model would spend capacity learning to reproduce the prompt.
\end{enumerate}

\section*{Chapter \ref{ch:measuring}: measurement}
\begin{enumerate}
  \item Accuracy can stay acceptable while one class is never predicted; macro-F1 averages per-class F1, so a missing class pulls it down.
  \item With 95\,\% confidence the true accuracy lies between 43.3\,\% and 68.6\,\%; the estimate is uncertain by about $\pm13$ points.
  \item Only the patches where the two disagree. Qwen got 31 patches right that the baseline got wrong, and missed all 19 increase patches the baseline got right; 31 against 19 is not unusual enough at this size ($p = 0.12$).
  \item Groundedness: the explanation agrees with the input. Faithfulness: it reflects the model's actual reasoning. The project measures groundedness (evidence consistency).
  \item A true small change below the dead-band (for example EVI $-0.006$ called a decrease) is marked wrong; a negated sentence (``did not decrease'') is read as a decrease claim.
\end{enumerate}

\section*{Chapter \ref{ch:trackA}: Track A}
\begin{enumerate}
  \item The card's trend word (for example \texttt{strong\_decline}) and the gold label \texttt{candidate\_class}.
  \item Qwen did not copy the label token, but the trend words it followed appear only in the same summary line; without the line it falls to chance.
  \item ``Stable'' for every patch in the compact run, which scores 18/55 = 32.7\,\% (the share of stable patches); the \texttt{evidence\_only} run also scored 32.7\,\%.
  \item No: it is not significantly above the 34.5\,\% baseline ($p = 0.44$), and it never answers ``stable''; its interval (31.4--56.7\,\%) includes chance.
  \item PaliGemma is gated and could not be downloaded without a token; BLIP-2 only echoed its prompt, so after the parser fix none of its answers contained a class.
  \item The information to classify is on the card as one number; the models do not use it.
\end{enumerate}

\section*{Chapter \ref{ch:trackB}: Track B}
\begin{enumerate}
  \item Job 29575 used 16 class-balanced patches (6/5/5) and the label-free card at training and test; job 2722446 used the first 16 patches by ID (9/6/1) and the full card.
  \item ``Stable'' for 53 of 55 patches; 36.4\,\% is the share it gets by mostly answering one class, no better than always answering ``increase'' ($p = 1.0$).
  \item It was obtained with the full card, which states the label, from an unbalanced sample; the leak-free rerun gives 36.4\,\%.
  \item With 256 patches and the full card the model learned to read the label from the card's last line (100\,\%); that measures copying.
  \item A balanced, label-free LoRA on 48 or all 256 training patches, with early stopping on the development set.
\end{enumerate}

\section*{Chapter \ref{ch:trackC}: Track C}
\begin{enumerate}
  \item A compact table of numbers with no class or interpretive words, so that the label is not in the input and the model must work from the numbers.
  \item PaliGemma gave classes without reasons; BLIP-2 repeated its prompt; LLaVA looped and contradicted itself; Qwen (3 patches) blamed a rainfall increase when rain fell.
  \item \texttt{ALLOWED\_FACTS}: numbered true sentences the model must copy verbatim. \texttt{CLASS\_RULE}: one line stating the class from the NDVI change with $\pm0.02$ cut-offs.
  \item 35 of Qwen's 38 and 35 of LLaVA's 37 flagged claims describe a true, tiny change that the scorer counts as stable because it is below the dead-band.
  \item Its class matched the stated rule on 53 of 55 patches, and the rule alone scores 63.6\,\%: it followed instructions.
  \item A 2B teacher passed only 28 of 256 justifications; the 16 used were all \texttt{decline}, and 9 of them were cut-off raw outputs; one epoch on 16 examples. (Any three.)
  \item The input (label-free table versus full card), the prompt (JSON explanation versus one word) and the training target all differ.
\end{enumerate}

\section*{Chapter \ref{ch:meaning}: what it shows}
\begin{enumerate}
  \item Classification: with the label line removed, no small VLM classifies better than guessing, although one number on the card gives 94.5\,\%. Explanations: consistent with the evidence only under a structured cite-only prompt, and then mostly restatements.
  \item H1, H2 (no images-only condition) and H4 (no usable confidence values).
  \item Data: relative labels, 19 human checks, one place and year pair, small changes and a dry year, padding bias. Models and evaluation: the summary line, unequal inputs in the first run, tiny fine-tunes, small test set, scorer limits (and others in Section~\ref{sec:limits}).
  \item Running zero-shot Qwen with the same free-form explanation prompt on all 55 test patches: it would show whether the decline/increase separation comes from the prompt or from the LoRA.
\end{enumerate}

\section*{Chapter \ref{ch:figures}: figures}
\begin{enumerate}
  \item \file{result2.png} (wrong encoder; a table as an image), \file{summary_scores.png} (outdated, 3-patch Qwen bars), \file{result1.png} and \file{ecoexplain_vlm_pipeline.png} (``LLaVA-1.5'', outdated statements).
  \item It says all four models received ``T0/T1 RGB + card'' and implies all four were run with and without the line; only Qwen and LLaVA saw the full card, and the \texttt{evidence\_only} re-run changed only LLaVA's input.
  \item The whole of India, not trimmed at the top, using the official boundary.
\end{enumerate}

\section*{Chapter \ref{ch:draft}: your draft}
\begin{enumerate}
  \item $56.4 - 34.5 = 21.9$ is a difference of 21.9 \emph{percentage points}; ``21.9\,\% higher'' would mean $1.219 \times 34.5 = 42.1\,\%$.
  \item PaliGemma received the card's NDVI line, not the ``summary line'' (which is the line containing the label); the sentence says the opposite of what happened.
  \item The final dataset's text is a rule-based narrative computed from the rasters (\texttt{raster\_analysis\_v2}); only the old version~1 was template text.
  \item Any two of: citations removed; a fourth heading level; missing acknowledgements, map and equations; Limitations placed after the Conclusions.
\end{enumerate}

\section*{Chapter \ref{ch:writing}: writing plan}
\begin{enumerate}
  \item Focus, motivation, gap and objectives; the gap is supported by \citet{jacovi2020}, \citet{turpin2023}, \citet{li2023pope} and \citet{zhao2021calibrate}.
  \item NDVI \citet{rouse1974,tucker1979}; EVI \citet{huete2002}; NDMI \citet{gao1996}; NBR \citet{keybenson2006}.
  \item Repetition adds words but no information, and two copies of a number can drift apart when one is corrected.
\end{enumerate}
